\chapter{Cosmological Spacetimes And Particle Creation}
\label{ch:curved-cosmological-spacetimes-particle-creation}

\ConventionsLorentzian

This chapter studies QFT on time-dependent globally hyperbolic backgrounds.
Particle creation is formulated through the relation between two choices of
positive-frequency decomposition or through detector response, with the state
and asymptotic conditions stated explicitly.

\section{Robertson-Walker Backgrounds}

Let
\[
  ds^2=-dt^2+a(t)^2 d\mathbf x^2
\]
be a spatially flat Robertson-Walker metric in \(d\) spacetime dimensions.
In conformal time \(d\eta=dt/a(t)\),
\[
  ds^2=a(\eta)^2(-d\eta^2+d\mathbf x^2).
\]
For a real scalar field,
\[
  (\Box_g-m^2-\xi R)\phi=0 .
\]
After Fourier transformation and the rescaling
\[
  \chi_{\mathbf k}(\eta)=a(\eta)^{(d-2)/2}\phi_{\mathbf k}(\eta),
\]
the mode equation has the form
\[
  \chi_{\mathbf k}''+\Omega_k(\eta)^2\chi_{\mathbf k}=0,
\]
where \(\Omega_k(\eta)^2\) is determined by \(k^2\), \(m^2a^2\), curvature
coupling, and derivatives of \(a\).

\section{Symplectic Space And Complex Structures}

Let \(\mathcal S\) be the real vector space of smooth solutions with compact
Cauchy data.  The conserved symplectic form is
\[
  \sigma(\phi_1,\phi_2)
  =
  \int_{\Sigma} d\Sigma^\mu\,
  (\phi_1\nabla_\mu\phi_2-\phi_2\nabla_\mu\phi_1).
\]
A quasifree Fock representation is specified by a complex structure
\[
  J:\mathcal S\to\mathcal S,\qquad J^2=-1,\qquad
  \sigma(Ju,Jv)=\sigma(u,v),
\]
such that \(\sigma(u,Ju)\) is positive for nonzero \(u\).  The corresponding
one-particle inner product is
\[
  \langle u,v\rangle_J=\sigma(u,Jv)+\ii\sigma(u,v).
\]

\section{Bogoliubov Coefficients}

Suppose the scale factor approaches static limits in the far past and far
future so that positive-frequency modes \(u_k^{\mathrm{in}}\) and
\(u_k^{\mathrm{out}}\) are defined.  They are related by
\[
  u_k^{\mathrm{in}}
  =
  \alpha_k u_k^{\mathrm{out}}
  +
  \beta_k \overline{u_k^{\mathrm{out}}},
\]
with
\[
  |\alpha_k|^2-|\beta_k|^2=1
\]
for each real scalar mode.  If the state is the in-vacuum, the expected
number of out-particles in mode \(k\) is
\[
  N_k^{\mathrm{out}}=|\beta_k|^2 .
\]
This number is tied to the two asymptotic complex structures; it is not an
additional local observable independent of that choice.

\section{Adiabatic States}

When no static asymptotic regions are available, adiabatic states provide a
local high-momentum condition.  Choose functions \(W_k(\eta)\) and write
\[
  \chi_k(\eta)=\frac{1}{\sqrt{2W_k(\eta)}}
  \exp\!\left(-\ii\int^\eta W_k(\eta')\,d\eta'\right).
\]
Substitution into the mode equation gives
\[
  W_k^2
  =
  \Omega_k^2
  -\frac12\frac{W_k''}{W_k}
  +\frac34\left(\frac{W_k'}{W_k}\right)^2 .
\]
Iterating this equation defines adiabatic order.  The ultraviolet subtraction
of the stress tensor in Robertson-Walker backgrounds is fixed by the same
large-\(k\) expansion as the Hadamard point-splitting prescription.

\section{Detector Response}

An Unruh-DeWitt detector on a timelike worldline \(x(\tau)\) with energy gap
\(E\) has transition probability, to leading order in the detector coupling,
\[
  P(E)
  =
  \int d\tau\,d\tau'\,
  e^{-\ii E(\tau-\tau')}
  \chi(\tau)\chi(\tau')
  W(x(\tau),x(\tau')),
\]
where \(\chi\) is a switching function and \(W\) is the two-point function of
the state.  Detector response is a local operational probe of the state along
the worldline and remains meaningful when a global particle decomposition is
not selected.

\section{de Sitter Example}

In \(d\)-dimensional de Sitter space with conformal time \(\eta<0\),
\[
  ds^2=\frac{1}{H^2\eta^2}(-d\eta^2+d\mathbf x^2).
\]
The Bunch-Davies quasifree state is selected by requiring that high-frequency
modes behave as positive-frequency Minkowski modes as \(\eta\to-\infty\):
\[
  \chi_k(\eta)\sim \frac{e^{-\ii k\eta}}{\sqrt{2k}} .
\]
For massive scalar fields this state is Hadamard and invariant under the
connected de Sitter group.  Other de Sitter-invariant two-point functions
must be checked against the Hadamard condition before they are accepted as
states of the local QFT.

\section{Backreaction Boundary}

The expectation value \(\langle T_{\mu\nu}\rangle\) in a time-dependent state
is defined by the renormalized stress-tensor prescription of
Chapter~\ref{ch:point-splitting-stress-tensor-renormalization}.  Semiclassical
Einstein equations,
\[
  G_{\mu\nu}+\Lambda g_{\mu\nu}
  =
  8\pi G\,\langle T_{\mu\nu}\rangle_{\mathrm{ren}},
\]
are equations for a classical metric coupled to a quantum state.  Solving
them requires state selection, renormalization constants, and control of the
validity of the fixed-background QFT approximation.

\begin{openproblem}[Interacting cosmological QFT]
\label{op:interacting-cosmological-qft}
Construct interacting QFT on cosmological spacetimes with Hadamard states,
renormalized stress tensor, detector response, particle-production diagnostics
where asymptotic regions exist, and controlled semiclassical backreaction.
\end{openproblem}
